\subsection{Automatic per-net release: completed native validation}
The independent automatic release build and validation have now completed. Eight tiny native children confirm unchanged default OFF/ON/failure results, guided-failure recovery with exactly the free native trees, distinct feedback request identity and a following default request unaffected in the same parent. On the actual BNN eight-cut scope, all eight domains are initially active; the native worker itself releases31767/pin1 after the guided search fails. Earlier restored trees and their occupancy are retained. Seven networks stay guided. The actual serialized candidate route and every fixed endpoint setup/hold component exactly match the prior selective-domain reference, with native FULL, physical/resource/domain audits, candidate fresh STA and final input/cache/context seals passed. The result remains13.228145689936355 ns,2069104 wire,72 SB and5071 CB. This is automatic recovery of the known resource point, not an additional independent frontier point.

The whole transaction is accepted under the registered resource-under-timing rule. Strict CPD rejects it, while endpoint-vector acceptance improves nine setup components with zero setup regressions and no increased hold violation. The worst endpoint identity is unchanged. In a fixed diagnostic source band of setup slack within0.05 times source CPD of the worst source slack, all23 components stay exact. Thus this step reduces actual resource occupancy and improves noncritical components without reducing the currently near-critical bottlenecks. Equality of endpoint values alone is not used to claim identical internal winning paths or future CPD benefit. The band is descriptive; it is not an acceptance tolerance and cannot accumulate CPD degradation.

The complete imported-state validation workflow is287.624 s, including source/native worker113.480 s, source materialization21.221 s, actual transaction audit41.032 s, domain audit10.984 s, candidate FULL/STA/materialization65.870 s, final context seal14.856 s and20.182 s of registration/context/miscellaneous work. Native repair core is6.693 s, nested within the worker rather than additive. Source/candidate native children account for43.712/37.073 s within these stages. No cold initial route is included. This run uses an exact sealed physical reference for the matching candidate, so neither its comparison with339.716 s of the earlier diagnostic nor its6.693 s core is advertised as a whole-algorithm speedup.

A subsequent independent cohort uses K8, the same0.1\% root-wire envelope and exact root-CPD cap from pristine strong native A0 starts for both eltwise and BNN. It compares native-only trajectories with paired automatic-feedback trajectories; each branch reselects after its own commits. No release network or accepted repair prefix is injected. These are prospective development probes; their still-running states are not included as completed quality points or timing results in this table. The two-step local probe is not an overall research stopping condition.
